\subsection{Chữ ký mù tập thể}
\label{sec:chu-ky-mu-tap-the}

Trong nhiều ứng dụng thực tế hiện đại, một thông điệp không chỉ cần được chứng thực bởi một cá nhân mà phải được chứng thực đồng thời bởi nhiều thực thể trong một tổ chức -- chẳng hạn ban tổ chức bầu cử gồm nhiều thành viên cùng đồng ý cấp phiếu cho cử tri, hoặc hội đồng thanh toán gồm nhiều ngân hàng cùng chứng thực một đồng tiền điện tử. Lược đồ chữ ký số tập thể (multi-signature) đáp ứng nhu cầu này nhưng chưa giải quyết được yêu cầu về tính riêng tư của người yêu cầu. Ngược lại, lược đồ chữ ký mù (mục 1.6) đảm bảo tính riêng tư nhưng chỉ có một người ký duy nhất. \textit{Chữ ký mù tập thể} (blind multi-signature -- BMS) ra đời như một sự kết hợp giữa hai loại chữ ký đặc biệt trên: cho phép nhiều người ký cùng ký lên một thông điệp đã được làm mù trong khi vẫn giữ kín nội dung thông điệp đối với mọi người ký. Đây chính là công cụ mật mã cốt lõi được sử dụng trong hệ thống bỏ phiếu điện tử blockchain được đề xuất trong đồ án này.

\subsubsection{Khái niệm chữ ký mù tập thể}

Lược đồ chữ ký mù tập thể cho phép nhiều người ký cùng ký trên một thông điệp đã được làm mù, sao cho không một người ký nào -- kể cả khi tất cả thoả hiệp với nhau -- có thể nhận biết được nội dung thực của thông điệp hay liên kết được chữ ký công bố sau này với phiên ký cụ thể trong quá khứ. Một cách hình thức, trong lược đồ chữ ký mù tập thể, người yêu cầu $A$ là chủ thể muốn nhận một chữ ký từ một \textit{tập thể người ký} $\mathbf{U} = \{U_1, U_2, \ldots, U_n\}$ trên thông điệp $M$ của mình; mỗi người ký $U_i$ sẽ không biết bất kỳ thông tin nào về mối quan hệ giữa thông điệp đã bị làm mù, thông điệp chưa bị làm mù và các tham số của chữ ký. Hệ quả là, ngay cả khi tất cả $n$ người ký thoả hiệp lại với nhau sau khi chữ ký đã được công bố, họ vẫn không thể nhận dạng được chữ ký đó thuộc về phiên ký nào trong các phiên ký đã thực hiện trước đây.

Công trình đầu tiên về chữ ký số tập thể mù được Horster và các cộng sự công bố vào năm 1995 trong bài báo ``\textit{Meta-Multisignature Schemes Based on the Discrete Logarithm Problem}'', đề xuất ứng dụng cho hệ thống bầu cử điện tử dựa trên hệ mật DLP. Một cách tiếp cận đơn giản nhất là mỗi người ký tự ký lên thông điệp đã làm mù bằng lược đồ chữ ký mù riêng lẻ rồi ghép các chữ ký lại với nhau; tuy nhiên, cách tiếp cận này có nhược điểm là kích thước chữ ký và thời gian xác minh tăng tuyến tính theo số lượng người ký, không phù hợp với các hệ thống quy mô lớn hoặc môi trường có tài nguyên hạn chế như blockchain. Do đó, các nghiên cứu sau này đều tập trung xây dựng các lược đồ chữ ký mù tập thể \textit{tích hợp}, trong đó chữ ký tập thể có kích thước cố định không phụ thuộc vào số lượng người ký và có thể được xác minh trong một bước duy nhất bằng khoá công khai tập thể.

\subsubsection{Mô hình tổng quát của chữ ký mù tập thể}

Một lược đồ chữ ký mù tập thể có thể được xem như là sự tổng quát hóa đồng thời của cả hai mô hình: mô hình chữ ký mù (mục 1.6) và mô hình chữ ký tập thể. Về mặt cấu trúc, một lược đồ chữ ký mù tập thể bao gồm bốn thuật toán chính:

\begin{itemize}
    \item \textbf{Thuật toán sinh khóa (Gen):} thường được thực hiện một lần ban đầu cho tất cả $n$ thành viên trong tập thể người ký. Mỗi thành viên $U_i$ nhận đầu vào là thông tin về nhóm $\mathbf{U}$ và sinh ra cặp khóa bí mật -- khóa công khai $(sk_i, pk_i)$ của riêng mình. Từ tập các khóa công khai riêng lẻ, hệ thống tính toán \textit{khóa công khai tập thể} $pk_{\mathbf{U}}$ làm khóa xác minh chung.

    \item \textbf{Giao thức làm mù (Blind)} được người yêu cầu $A$ thực hiện cục bộ trên thông điệp $M$ với một bộ tham số làm mù ngẫu nhiên bí mật, tạo ra phiên bản $M'$ đã được làm mù để gửi cho tập thể người ký.

    \item \textbf{Giao thức ký tập thể (BlindSig)}: là một giao thức tương tác nhiều bên giữa người yêu cầu $A$ và các thành viên trong tập thể $\mathbf{U}$. Trong giao thức này, mỗi người ký đóng góp một phần chữ ký riêng từ khoá bí mật của mình; các phần chữ ký riêng được tổng hợp thành chữ ký tập thể bởi $A$ hoặc bởi một bên thứ ba tin cậy (Trusted Third Party -- TTP).

    \item \textbf{Thuật toán giải mù (Unblind)} và \textbf{thuật toán xác minh (Ver)}: $A$ thực hiện giải mù để thu được chữ ký hợp lệ trên thông điệp gốc $M$. Bất kỳ thực thể nào trong hệ thống cũng có thể xác minh tính hợp lệ của cặp $(M, \sigma)$ bằng khoá công khai tập thể $pk_{\mathbf{U}}$ và bộ tham số công khai của lược đồ.
\end{itemize}

Tiến trình tổng quát của một lược đồ chữ ký mù tập thể được mô tả khái quát như sau: người yêu cầu $A$ cần tập thể $\mathbf{U}$ ký cho thông điệp $M$, nhưng $A$ không đưa $M$ cho $\mathbf{U}$ ký mà làm mù $M$ thành $M'$ trước khi đưa $M'$ cho $\mathbf{U}$. Tập thể $\mathbf{U}$ thực hiện giao thức ký tập thể trên $M'$ và trả về phần ký kết quả $\sigma'$ cho $A$. Sau khi nhận được $\sigma'$, $A$ giải mù để thu được chữ ký $\sigma$ trên $M$. Như vậy $A$ thu được chữ ký hợp lệ của tập thể $\mathbf{U}$ trên thông điệp $M$ mà các thành viên của $\mathbf{U}$ không biết bất kỳ thông tin nào về $M$.

\subsubsection{Vai trò của các thực thể tham gia}

Trong một lược đồ chữ ký mù tập thể, có bốn vai trò chính có thể được phân biệt:

\begin{itemize}
    \item \textbf{Người yêu cầu} ($A$ -- Requester) là chủ thể sở hữu thông điệp $M$ và mong muốn nhận được chữ ký của tập thể $\mathbf{U}$ trên $M$. $A$ là thực thể duy nhất biết được nội dung thực của $M$ và nắm giữ các tham số làm mù bí mật $(\alpha, \beta, \ldots)$ được sử dụng trong các pha làm mù và giải mù. Sau khi quá trình ký hoàn tất, $A$ có thể công bố cặp $(M, \sigma)$ một cách nặc danh mà không lo lộ danh tính của mình. Trong bối cảnh bỏ phiếu điện tử, $A$ chính là cử tri.

    \item \textbf{Tập thể người ký} ($\mathbf{U} = \{U_1, \ldots, U_n\}$ -- Signing Group) là nhóm các thực thể có thẩm quyền chứng thực và đóng vai trò ký xác nhận quyền của $A$. Mỗi thành viên $U_i$ nắm giữ một khoá bí mật riêng $sk_i = d_i$ và đóng góp một phần chữ ký vào chữ ký tập thể cuối cùng. Đặc điểm quan trọng là không một $U_i$ nào (kể cả khi tất cả thoả hiệp) có thể thu được thông tin về $M$ hay liên kết chữ ký đã công bố với phiên ký cụ thể nào. Trong hệ thống bỏ phiếu điện tử, $\mathbf{U}$ tương ứng với ban tổ chức bầu cử gồm nhiều thành viên.

    \item \textbf{Bên thứ ba tin cậy} (TTP -- Trusted Third Party) là thực thể tùy chọn có nhiệm vụ điều phối giao thức ký tập thể: thu thập các phần chữ ký riêng từ các thành viên, tổng hợp lại thành chữ ký tập thể và trả về cho người yêu cầu. Trong một số phiên bản, vai trò này có thể được phân tán hoặc gán cho chính người yêu cầu $A$ để giảm yêu cầu về tin cậy hệ thống.

    \item \textbf{Người xác minh} (Verifier) là bất kỳ thực thể nào trong hệ thống có thể truy cập đến bộ tham số công khai và khoá công khai tập thể $pk_{\mathbf{U}}$. Người xác minh không tham gia vào quá trình ký mà chỉ thực hiện kiểm tra tính hợp lệ của cặp $(M, \sigma)$ được công bố cuối cùng.
\end{itemize}

\subsubsection{Quy trình tạo chữ ký mù tập thể}

Quy trình tạo chữ ký mù tập thể bao gồm năm pha tuần tự được mô tả như sau:

\begin{itemize}
    \item \textbf{Pha 1 -- Khởi tạo và sinh khóa:} Mỗi thành viên $U_i \in \mathbf{U}$ độc lập sinh khóa bí mật $d_i$ và tính khóa công khai $P_i$. Khóa công khai tập thể $P$ được tính toán bằng phép tổng hợp các khóa công khai riêng lẻ và được công bố cùng với danh sách định danh của các thành viên trong $\mathbf{U}$.

    \item \textbf{Pha 2 -- Sinh cam kết tập thể:} Mỗi thành viên $U_i$ độc lập chọn giá trị ngẫu nhiên $k_i$ và tính điểm cam kết $C_i$. Các điểm cam kết được tập hợp và tổng hợp thành điểm cam kết tập thể $C$ rồi gửi cho người yêu cầu $A$.

    \item \textbf{Pha 3 -- Làm mù:} Người yêu cầu $A$ sử dụng các tham số làm mù bí mật $(\alpha, \beta)$ được chọn ngẫu nhiên để biến đổi điểm cam kết tập thể $C$ thành phiên bản đã làm mù $C'$, sau đó tính giá trị thử thách đã làm mù $e$ dựa trên $C'$ và thông điệp $M$, rồi gửi $e$ về cho tập thể người ký.

    \item \textbf{Pha 4 -- Ký tập thể:} Mỗi thành viên $U_i$ nhận thử thách $e$ và sử dụng khóa bí mật $d_i$ cùng giá trị ngẫu nhiên $k_i$ để tính phần chữ ký riêng $s_i$. Các phần chữ ký riêng được tổng hợp thành phần chữ ký tập thể đã làm mù $s'$ và gửi về cho $A$.

    \item \textbf{Pha 5 -- Giải mù và xác minh:} Người yêu cầu $A$ áp dụng các tham số làm mù để giải mù $s'$, thu được chữ ký hoàn chỉnh $\sigma$ trên thông điệp gốc $M$. Bất kỳ thực thể nào cũng có thể xác minh $(M, \sigma)$ bằng khóa công khai tập thể $P$ thông qua thuật toán xác minh chuẩn của lược đồ.
\end{itemize}

\subsubsection{Lược đồ chữ ký mù tập thể dựa trên EC-Schnorr}

Trên cơ sở lược đồ EC-Schnorr (mục 1.5) và mô hình chữ ký mù (mục 1.6), lược đồ chữ ký mù tập thể dựa trên EC-Schnorr được xây dựng cụ thể như sau. Bộ tham số miền gồm đường cong elliptic $E$ trên trường $\mathbb{F}_p$, điểm gốc $G$ cấp $q$, hàm băm $H: \{0,1\}^* \rightarrow \mathbb{Z}_q$, tập thể $n$ người ký $\mathbf{U} = \{U_1, \ldots, U_n\}$ và người yêu cầu $A$.

\begin{itemize}
    \item \textbf{Pha sinh khóa:} Mỗi thành viên $U_i$ chọn ngẫu nhiên $d_i \in (1, q-1)$ làm khóa bí mật và tính khóa công khai riêng $P_i = d_i \times G$. Khóa công khai tập thể được tính bằng tổng các khóa công khai riêng:
    \begin{equation}
        P = \sum_{i=1}^{n} P_i = \left( \sum_{i=1}^{n} d_i \right) \times G = D \times G \bmod p
        \label{eq:bms-pubkey}
    \end{equation}
    trong đó $D = \sum_{i=1}^{n} d_i \bmod q$ là khoá bí mật tập thể (đại lượng này không được tính toán tường minh và không một thành viên đơn lẻ nào biết được).

    \item \textbf{Pha sinh cam kết:} Mỗi thành viên $U_i$ chọn ngẫu nhiên $k_i \in (1, q-1)$ và tính điểm cam kết riêng:
    \begin{equation}
        C_i = k_i \times G \bmod p
        \label{eq:bms-commit-i}
    \end{equation}
    Các điểm $C_i$ được gửi tới $A$ (hoặc TTP), nơi tổng hợp thành điểm cam kết tập thể:
    \begin{equation}
        C = \sum_{i=1}^{n} C_i = \left( \sum_{i=1}^{n} k_i \right) \times G = K \times G \bmod p
        \label{eq:bms-commit}
    \end{equation}
    với $K = \sum_{i=1}^{n} k_i \bmod q$.

    \item \textbf{Pha làm mù:} Người yêu cầu $A$ chọn ngẫu nhiên hai tham số làm mù bí mật $\alpha, \beta \in \mathbb{Z}_q$ và tính điểm cam kết đã làm mù:
    \begin{equation}
        C' = C + \alpha \times G + \beta \times P \bmod p
        \label{eq:bms-blind-c}
    \end{equation}
    Sau đó $A$ tính giá trị $r'$ -- thành phần thứ nhất của chữ ký cuối cùng -- và giá trị $e$ -- thử thách đã làm mù sẽ được gửi về cho người ký:
    \begin{equation}
        r' = H(M, x_{C'}) \bmod q, \qquad e = (r' - \beta) \bmod q
        \label{eq:bms-blind-challenge}
    \end{equation}
    trong đó $x_{C'}$ là hoành độ của điểm $C'$. Giá trị $e$ được gửi tới tất cả các thành viên trong $\mathbf{U}$.

    \item \textbf{Pha ký tập thể:} Mỗi thành viên $U_i$ nhận $e$ và tính phần chữ ký riêng:
    \begin{equation}
        s_i = (k_i - e \cdot d_i) \bmod q
        \label{eq:bms-sign-i}
    \end{equation}
    Sau đó gửi $s_i$ về cho $A$. $A$ tổng hợp các phần chữ ký riêng thành phần chữ ký tập thể đã làm mù:
    \begin{equation}
        s'' = \sum_{i=1}^{n} s_i \bmod q
        \label{eq:bms-sign-aggregate}
    \end{equation}

    \item \textbf{Pha giải mù:} Người yêu cầu $A$ áp dụng tham số $\alpha$ để giải mù, thu được thành phần thứ hai của chữ ký:
    \begin{equation}
        s' = (s'' + \alpha) \bmod q
        \label{eq:bms-unblind}
    \end{equation}
    Cặp $(r', s')$ chính là chữ ký mù tập thể trên thông điệp $M$.

    \item \textbf{Pha xác minh:} Để xác minh tính hợp lệ của chữ ký, người xác minh tính:
    \begin{equation}
        C^* = s' \times G + r' \times P \bmod p, \qquad r^* = H(M, x_{C^*}) \bmod q
        \label{eq:bms-verify}
    \end{equation}
    Nếu $r^* = r'$ thì chữ ký được chấp nhận, ngược lại bị từ chối.
\end{itemize}

\subsubsection{Phân tích tính đúng đắn}

\textit{Định lý} (Tính đúng đắn): Với bất kỳ thông điệp $M$, bất kỳ tập khóa $\{d_i\}$ và các giá trị ngẫu nhiên hợp lệ $\{k_i\}, \alpha, \beta$, chữ ký mù tập thể $(r', s')$ được tạo bởi giao thức trên luôn được thuật toán xác minh chấp nhận.

\textit{Chứng minh}: Thay thế \eqref{eq:bms-unblind}, \eqref{eq:bms-sign-aggregate}, \eqref{eq:bms-sign-i} và \eqref{eq:bms-blind-challenge} vào vế phải của \eqref{eq:bms-verify}, ta có:
\begin{align}
    C^* &= s' \times G + r' \times P \notag \\
        &= (s'' + \alpha) \times G + r' \times P \notag \\
        &= \left( \sum_{i=1}^{n} (k_i - e \cdot d_i) + \alpha \right) \times G + r' \times P \notag \\
        &= \left( \sum_{i=1}^{n} k_i - e \cdot \sum_{i=1}^{n} d_i + \alpha \right) \times G + r' \times P \notag \\
        &= (K - e D + \alpha) \times G + r' \times P \notag \\
        &= K \times G - e \times (D \times G) + \alpha \times G + r' \times P \notag \\
        &= C - e \times P + \alpha \times G + r' \times P \notag \\
        &= C + \alpha \times G + (r' - e) \times P
        \label{eq:bms-proof-1}
\end{align}
Theo định nghĩa \eqref{eq:bms-blind-challenge}, ta có $e = r' - \beta$, do đó $r' - e = \beta$, suy ra:
\begin{equation}
    C^* = C + \alpha \times G + \beta \times P = C'
    \label{eq:bms-proof-2}
\end{equation}
Như vậy $x_{C^*} = x_{C'}$, dẫn đến:
\begin{equation}
    r^* = H(M, x_{C^*}) \bmod q = H(M, x_{C'}) \bmod q = r'
    \label{eq:bms-proof-3}
\end{equation}
Điều kiện xác minh $r^* = r'$ luôn được thoả mãn, do đó chữ ký hợp lệ luôn được chấp nhận. Tính đúng đắn của lược đồ đã được chứng minh. $\hfill \blacksquare$

\subsubsection{Phân tích tính an toàn}

Tính an toàn của lược đồ chữ ký mù tập thể dựa trên EC-Schnorr được đánh giá thông qua bốn khía cạnh tương ứng với các thuộc tính an toàn của chữ ký mù tập thể.

\begin{itemize}
    \item \textbf{Tính mù (Blindness):} Đối với mỗi người ký $U_i$, các giá trị quan sát được trong giao thức gồm điểm cam kết $C_i$ do chính mình tạo, thử thách $e$ và phần chữ ký $s_i$. Do $A$ chọn ngẫu nhiên hai tham số làm mù $\alpha, \beta \in \mathbb{Z}_q$ phân bố đều, ánh xạ $(r', e) \mapsto (\alpha, \beta)$ với $\beta = r' - e$ và $\alpha$ được dùng trong $C' = C + \alpha G + \beta P$ là song ánh đối với mỗi cặp $(C, e)$ cho trước. Do đó, từ góc nhìn của người ký, mọi cặp $(M, r', s')$ hợp lệ đều có cùng phân bố xác suất với bất kỳ phiên ký cụ thể nào trong quá khứ. Điều này đảm bảo tính mù theo nghĩa của Định nghĩa 1.1 -- xác suất phân biệt được hai phiên ký khác nhau là một hàm vô cùng nhỏ trong cấp $q$ của nhóm.

    \item \textbf{Tính không thể giả mạo (Unforgeability):} Tính không thể giả mạo của lược đồ được chứng minh trong mô hình ROM thông qua kỹ thuật quy giảm về bài toán ECDLP. Cụ thể, giả sử tồn tại một đối thủ $\mathcal{A}$ với thuật toán hiệu năng cao trong thời gian đa thức có thể giả mạo chữ ký mù tập thể với xác suất $\varepsilon$ không bỏ qua được; khi đó có thể xây dựng một thuật toán $\mathcal{B}$ giải bài toán ECDLP với xác suất $\varepsilon' = \left(1 - \frac{q_h(q_e + q_s)}{q}\right)\left(1 - \frac{1}{q}\right)\frac{\varepsilon}{q_h}$ thông qua kỹ thuật \textit{Forking Lemma} của Pointcheval và Stern, trong đó $(q_h, q_e, q_s)$ lần lượt là số truy vấn tới hàm băm, hàm trích xuất và hàm ký. Do bài toán ECDLP được giả thiết là khó với các tham số chuẩn, kết quả này mâu thuẫn với giả thiết tồn tại $\mathcal{A}$, do đó lược đồ là an toàn EUF-CMA.

    \item \textbf{Tính không thể liên kết (Unlinkability):} Do tính mù được đảm bảo và các tham số làm mù $(\alpha, \beta)$ được chọn độc lập, ngẫu nhiên cho mỗi phiên ký, không có bất kỳ thông tin tương quan thống kê nào giữa các giá trị quan sát được trong phiên ký $(C, e, \{s_i\})$ và chữ ký cuối cùng $(M, r', s')$. Hệ quả là, ngay cả khi tất cả $n$ người ký thoả hiệp lại với nhau sau khi chữ ký được công bố, xác suất xác định đúng phiên ký tương ứng với một chữ ký công bố cụ thể chỉ bằng $1/N$ với $N$ là số phiên ký đã thực hiện -- tương đương với việc đoán ngẫu nhiên.

    \item \textbf{An toàn của khóa bí mật và yêu cầu kỹ thuật:} Khóa bí mật tập thể $D = \sum d_i$ chưa bao giờ được tính toán tường minh và không một thành viên đơn lẻ nào biết được; việc khôi phục $D$ từ $P = D \times G$ tương đương trực tiếp với bài toán ECDLP. Tương tự, mỗi giá trị ngẫu nhiên $k_i$ phải được sinh hoàn toàn ngẫu nhiên, giữ bí mật và không được tái sử dụng giữa các phiên ký, bởi nếu hai phần chữ ký $s_i$ trên hai thử thách $e$ khác nhau cùng dùng một $k_i$ thì khoá bí mật $d_i$ có thể bị khôi phục ngay lập tức. Ngoài ra, các thành viên trong tập thể cần một kênh truyền tin xác thực để chống lại các tấn công \textit{rogue-key} trong đó một người ký gian lận chọn khoá công khai phụ thuộc vào khoá công khai của những người khác nhằm chiếm đoạt chữ ký.
\end{itemize}

Với các phân tích trên, lược đồ chữ ký mù tập thể dựa trên EC-Schnorr đáp ứng đầy đủ các yêu cầu an toàn về tính mù, tính không thể giả mạo và tính không thể liên kết theo các định nghĩa 1.1 và 1.2 (mục 1.6.3), đồng thời kế thừa toàn bộ ưu điểm về hiệu năng và kích thước chữ ký của lược đồ EC-Schnorr. Đặc biệt, kích thước chữ ký tập thể cố định bằng $2 \lceil \log_2 q \rceil$ bit -- không phụ thuộc vào số lượng người ký $n$ -- và quá trình xác minh chỉ cần một phép kiểm tra duy nhất bằng khoá công khai tập thể $P$. Đây chính là những đặc tính khiến lược đồ này trở thành lựa chọn lý tưởng cho hệ thống bỏ phiếu điện tử dựa trên blockchain Hyperledger Fabric được đề xuất ở các chương tiếp theo, nơi yêu cầu đồng thời về tính nặc danh của cử tri, tính xác thực của ban tổ chức và hiệu năng giao dịch cao trên môi trường chuỗi khối.
